\usepackage[acronym,toc,nonumberlist]{glossaries}

% Separate glossary for mathematical and physical symbols.
\newglossary[slg]{symbols}{syi}{syg}{List of Symbols}

% No-index generation is convenient for a self-contained thesis build.
\makenoidxglossaries

% Keep glossary links in the surrounding text colour.
\renewcommand{\glstextformat}[1]{#1}

% ============================================================================
% ACRONYMS
% ============================================================================

\newacronym{ae}{AE}{autoencoder}
\newacronym{ai}{AI}{artificial intelligence}
\newacronym{bo}{BO}{Bayesian optimization}
\newacronym{cern}{CERN}{European Organization for Nuclear Research}
\newacronym{cnn}{CNN}{convolutional neural network}
\newacronym{cpu}{CPU}{central processing unit}
\newacronym{cuda}{CUDA}{Compute Unified Device Architecture}
\newacronym{cvae}{CVAE}{conditional variational autoencoder}
\newacronym{dl}{DL}{deep learning}
\newacronym{doi}{DOI}{digital object identifier}
\newacronym{elbo}{ELBO}{evidence lower bound}
\newacronym{fccee}{FCC-ee}{Future Circular Collider electron--positron collider}
\newacronym{fcc}{FCC}{Future Circular Collider}
\newacronym{fft}{FFT}{fast Fourier transform}
\newacronym{fno}{FNO}{Fourier Neural Operator}
\newacronym{ganil}{GANIL}{Grand Accélérateur National d'Ions Lourds}
\newacronym{gelu}{GELU}{Gaussian error linear unit}
\newacronym{gno}{GNO}{Graph Neural Operator}
\newacronym{gpu}{GPU}{graphics processing unit}
\newacronym{heb}{HEB}{High-Energy Booster}
\newacronym{ipac}{IPAC}{International Particle Accelerator Conference}
\newacronym{jacow}{JACoW}{Joint Accelerator Conferences Website}
\newacronym{kl}{KL}{Kullback--Leibler}
\newacronym{madx}{MAD-X}{Methodical Accelerator Design, version X}
\newacronym{mha}{MHA}{multi-head attention}
\newacronym{ml}{ML}{machine learning}
\newacronym{mlp}{MLP}{multilayer perceptron}
\newacronym{mse}{MSE}{mean squared error}
\newacronym{neg}{NEG}{non-evaporable getter}
\newacronym{ood}{OOD}{out of distribution}
\newacronym{pde}{PDE}{partial differential equation}
\newacronym{rf}{RF}{radio frequency}
\newacronym{rms}{RMS}{root mean square}
\newacronym{swd}{SWD}{sliced Wasserstein distance}
\newacronym{tf}{TF}{Transformer}
\newacronym{vae}{VAE}{variational autoencoder}
\newacronym{vfp}{VFP}{Vlasov--Fokker--Planck}
\newacronym{uq}{UQ}{uncertainty quantification}

% ============================================================================
% SOFTWARE, DATA, AND COMPUTATIONAL TOOLS
% ============================================================================

\newglossaryentry{xsuite}{
    name={Xsuite},
    sort={Xsuite},
    description={Open-source Python suite for multiparticle accelerator simulations, including tracking, optics manipulation, collective-effects interfaces, and accelerator-model workflows}
}

\newglossaryentry{xtrack}{
    name={Xtrack},
    sort={Xtrack},
    description={Xsuite component used to construct accelerator lines and track macroparticles through lattice elements}
}

\newglossaryentry{cpymad}{
    name={\texttt{cpymad}},
    sort={cpymad},
    description={Python interface to MAD-X used to load and manipulate accelerator lattice descriptions}
}

\newglossaryentry{pycolleff}{
    name={PyColleff},
    sort={PyColleff},
    description={Python module developed by LNLS for impedance modelling and collective-instability analysis in electron storage rings}
}

\newglossaryentry{pywit}{
    name={PyWIT},
    sort={PyWIT},
    description={Python Wake and Impedance Toolbox used to represent and manipulate accelerator wake and impedance models}
}

% \newglossaryentry{zenodo}{
%     name={Zenodo},
%     sort={Zenodo},
%     description={Open research repository used to archive datasets, software releases, and associated metadata with a persistent DOI}
% }

% \newglossaryentry{dataset-manifest}{
%     name={dataset manifest},
%     sort={dataset manifest},
%     description={Machine-readable record of dataset files, array shapes, parameter ranges, units, software versions, checksums, and train--validation--test indices}
% }

% \newglossaryentry{checkpoint}{
%     name={model checkpoint},
%     sort={model checkpoint},
%     description={Saved state of a trained model, usually including learned parameters and, when required, optimizer or preprocessing state}
% }

% ============================================================================
% ACCELERATOR AND BEAM-PHYSICS TERMS
% ============================================================================

\newglossaryentry{accelerator-lattice}{
    name={accelerator lattice},
    sort={accelerator lattice},
    description={Ordered sequence of magnets, radio-frequency systems, drifts, and other elements that determines charged-particle transport through an accelerator}
}

\newglossaryentry{beam-centroid}{
    name={beam centroid},
    sort={beam centroid},
    description={First statistical moment of a beam-coordinate distribution; for coordinate $\xi_i$, it is the ensemble average $\langle \xi_i\rangle$}
}

\newglossaryentry{beam-size}{
    name={beam size},
    sort={beam size},
    description={Transverse spatial extent of a beam, conventionally expressed through an RMS coordinate spread and related to emittance, optical beta function, dispersion, and momentum spread}
}

\newglossaryentry{betatron-motion}{
    name={betatron motion},
    sort={betatron motion},
    description={Oscillatory transverse motion of a particle around the reference orbit under the action of focusing elements}
}

\newglossaryentry{bunch}{
    name={bunch},
    sort={bunch},
    description={Localized ensemble of charged particles transported together through an accelerator}
}

\newglossaryentry{bunch-charge}{
    name={bunch charge},
    sort={bunch charge},
    description={Total electric charge carried by one particle bunch}
}

\newglossaryentry{collective-effects}{
    name={collective effects},
    sort={collective effects},
    description={Beam-dynamics phenomena in which particles interact through fields generated by the beam itself or by its electromagnetic interaction with the accelerator environment}
}

\newglossaryentry{coupled-bunch-instability}{
    name={coupled-bunch instability},
    sort={coupled-bunch instability},
    description={Coherent instability in which wake-mediated interactions couple the motion of different bunches circulating in an accelerator}
}

\newglossaryentry{dispersion}{
    name={dispersion},
    sort={dispersion},
    description={Dependence of the closed orbit or particle position on relative momentum deviation; it contributes to the measured transverse beam size}
}

\newglossaryentry{emittance}{
    name={emittance},
    sort={emittance},
    description={Phase-space area or volume characterizing the spread of a beam in a conjugate coordinate--momentum plane}
}

\newglossaryentry{haissinski}{
    name={Ha\"issinski distribution},
    sort={Haissinski distribution},
    description={Stationary longitudinal bunch distribution obtained from the self-consistent balance of RF confinement, wakefields, synchrotron-radiation damping, and quantum excitation below the microwave-instability threshold}
}

\newglossaryentry{impedance}{
    name={beam coupling impedance},
    sort={beam coupling impedance},
    description={Frequency-domain representation of the electromagnetic response of the accelerator environment to the beam}
}

\newglossaryentry{intrabeam-scattering}{
    name={intrabeam scattering},
    sort={intrabeam scattering},
    description={Multiple small-angle Coulomb scattering among particles within the same bunch, which can increase beam emittances and energy spread}
}

\newglossaryentry{line-density}{
    name={longitudinal line density},
    sort={longitudinal line density},
    description={Particle or charge density projected onto the longitudinal coordinate}
}

\newglossaryentry{liouville-theorem}{
    name={Liouville's theorem},
    sort={Liouville theorem},
    description={Statement that Hamiltonian phase-space flow preserves phase-space volume and transports densities incompressibly}
}

\newglossaryentry{macroparticle}{
    name={macroparticle},
    sort={macroparticle},
    description={Numerical simulation particle representing many physical particles while carrying effective charge and phase-space coordinates}
}

\newglossaryentry{microwave-instability}{
    name={microwave instability},
    sort={microwave instability},
    description={Single-bunch longitudinal instability associated with short-range wakefields and high-frequency impedance}
}

\newglossaryentry{panofsky-wenzel}{
    name={Panofsky--Wenzel theorem},
    sort={Panofsky Wenzel theorem},
    description={Relation connecting longitudinal and transverse wakefields through derivatives of the integrated electromagnetic force}
}

\newglossaryentry{phase-space}{
    name={phase space},
    sort={phase space},
    description={Space of particle positions and conjugate momenta; the beam state in this thesis is represented in six dimensions}
}

\newglossaryentry{phase-space-marginal}{
    name={phase-space marginal},
    sort={phase-space marginal},
    description={Lower-dimensional distribution obtained by integrating the full phase-space density over unobserved coordinates}
}

\newglossaryentry{quantile}{
    name={quantile},
    sort={quantile},
    description={Value below which a specified fraction of a probability distribution lies; useful for describing non-Gaussian beam tails and asymmetric distributions}
}

\newglossaryentry{resistive-wall}{
    name={resistive-wall impedance},
    sort={resistive wall impedance},
    description={Beam coupling impedance produced by finite electrical conductivity of the vacuum chamber walls}
}

\newglossaryentry{rms-spread}{
    name={coordinate-wise RMS spread},
    sort={coordinate-wise RMS spread},
    description={Square root of the second central moment of one beam coordinate; it remains defined for non-Gaussian distributions with finite variance but does not fully characterize their shape}
}

\newglossaryentry{self-consistent-force}{
    name={self-consistent force},
    sort={self consistent force},
    description={Force that depends on the instantaneous beam distribution and therefore couples the evolution of individual particles to the collective state}
}

\newglossaryentry{space-charge}{
    name={space charge},
    sort={space charge},
    description={Collective electromagnetic interaction arising directly from the charge and current distribution of the beam}
}

\newglossaryentry{synchrotron-radiation}{
    name={synchrotron radiation},
    sort={synchrotron radiation},
    description={Electromagnetic radiation emitted by relativistic charged particles undergoing transverse acceleration, producing damping and quantum excitation in electron rings}
}

\newglossaryentry{symplectic-map}{
    name={symplectic map},
    sort={symplectic map},
    description={Phase-space transformation that preserves the symplectic structure of Hamiltonian dynamics and therefore respects the associated geometric invariants}
}

\newglossaryentry{twiss-parameters}{
    name={Twiss parameters},
    sort={Twiss parameters},
    description={Optical functions $\alpha$, $\beta$, and $\gamma$ used to describe linear transverse beam envelopes and phase-space ellipses}
}

\newglossaryentry{vlasov-equation}{
    name={Vlasov equation},
    sort={Vlasov equation},
    description={Kinetic equation describing collisionless evolution of a phase-space density under external and self-consistent forces}
}

\newglossaryentry{vlasov-fokker-planck}{
    name={Vlasov--Fokker--Planck equation},
    sort={Vlasov Fokker Planck equation},
    description={Kinetic equation combining Vlasov transport with damping and diffusion terms}
}

\newglossaryentry{wake-function}{
    name={wake function},
    sort={wake function},
    description={Green-function response of the accelerator environment to a point-like source charge}
}

\newglossaryentry{wake-potential}{
    name={wake potential},
    sort={wake potential},
    description={Integrated electromagnetic effect experienced by a beam distribution, obtained by convolving the wake function with the source distribution}
}

% ============================================================================
% MATHEMATICAL AND STATISTICAL TERMS
% ============================================================================

\newglossaryentry{convolution}{
    name={convolution},
    sort={convolution},
    description={Integral operation combining a function with a translated kernel; in beam dynamics it relates a longitudinal density to its induced wake potential}
}

\newglossaryentry{empirical-measure}{
    name={empirical measure},
    sort={empirical measure},
    description={Discrete probability measure formed by assigning equal point mass to each simulated macroparticle}
}

\newglossaryentry{fourier-transform}{
    name={Fourier transform},
    sort={Fourier transform},
    description={Transformation between a function in physical space and its representation in frequency or wave-number space}
}

\newglossaryentry{kernel}{
    name={kernel},
    sort={kernel},
    description={Function defining weighted interactions or an integral operator between input and output locations}
}

\newglossaryentry{kl-divergence}{
    name={Kullback--Leibler divergence},
    sort={Kullback Leibler divergence},
    description={Asymmetric divergence measuring how one probability distribution differs from another; in a VAE it regularizes the approximate posterior toward the prior}
}

\newglossaryentry{monte-carlo}{
    name={Monte Carlo estimator},
    sort={Monte Carlo estimator},
    description={Approximation of an expectation or integral by an average over random samples}
}

\newglossaryentry{posterior}{
    name={posterior distribution},
    sort={posterior distribution},
    description={Probability distribution updated after conditioning on observed data or model inputs}
}

\newglossaryentry{prior}{
    name={prior distribution},
    sort={prior distribution},
    description={Probability distribution specified before conditioning on observations; the standard VAE prior is commonly a unit Gaussian}
}

\newglossaryentry{pushforward}{
    name={pushforward measure},
    sort={pushforward measure},
    description={Distribution obtained by transporting a measure through a deterministic map}
}

\newglossaryentry{wasserstein-distance}{
    name={Wasserstein distance},
    sort={Wasserstein distance},
    description={Optimal-transport distance between probability distributions that accounts for the cost of moving probability mass}
}

\newglossaryentry{sliced-wasserstein}{
    name={sliced Wasserstein distance},
    sort={sliced Wasserstein distance},
    description={Distributional distance obtained by averaging one-dimensional Wasserstein distances over projected directions}
}

% ============================================================================
% MACHINE-LEARNING AND SURROGATE-MODELLING TERMS
% ============================================================================

\newglossaryentry{attention}{
    name={attention},
    sort={attention},
    description={Neural-network mechanism that forms weighted combinations of value representations using learned similarities between queries and keys}
}

\newglossaryentry{autoencoder}{
    name={autoencoder},
    sort={autoencoder},
    description={Neural network trained to encode data into a lower-dimensional representation and decode that representation back to the input space}
}

\newglossaryentry{causal-attention}{
    name={causal attention},
    sort={causal attention},
    description={Attention mask that permits a token to receive information only from itself and preceding tokens in an ordered sequence}
}

\newglossaryentry{conditional-latent-uncertainty}{
    name={conditional latent uncertainty},
    sort={conditional latent uncertainty},
    description={Variability obtained by repeatedly sampling the conditional latent distribution while holding model parameters and conditioning variables fixed}
}

\newglossaryentry{calibration}{
    name={uncertainty calibration},
    sort={uncertainty calibration},
    description={Agreement between predicted uncertainty intervals or probabilities and their empirical coverage or observed error frequencies}
}

\newglossaryentry{decoder}{
    name={decoder},
    sort={decoder},
    description={Model component that maps a latent representation back to the data or observable space}
}

\newglossaryentry{digital-twin}{
    name={digital twin},
    sort={digital twin},
    description={Computational representation of a physical system that combines models and data to support prediction, monitoring, calibration, or control}
}

\newglossaryentry{discretization-invariance}{
    name={discretization invariance},
    sort={discretization invariance},
    description={Property of an operator-learning architecture whose parameters do not depend on one fixed input or output discretization and whose discrete realizations converge under refinement}
}

\newglossaryentry{encoder}{
    name={encoder},
    sort={encoder},
    description={Model component that maps high-dimensional observations into features or a latent representation}
}

\newglossaryentry{graph-message-passing}{
    name={graph message passing},
    sort={graph message passing},
    description={Update mechanism in which graph nodes aggregate learned messages from neighbouring nodes connected by edges}
}

\newglossaryentry{inductive-bias}{
    name={inductive bias},
    sort={inductive bias},
    description={Architectural assumption that favours particular solution structures, such as locality, causality, translation structure, or graph connectivity}
}

\newglossaryentry{latent-space}{
    name={latent space},
    sort={latent space},
    description={Learned lower-dimensional representation in which complex observations are encoded and manipulated}
}

\newglossaryentry{local-window-attention}{
    name={local-window attention},
    sort={local window attention},
    description={Attention mechanism restricted to a finite neighbourhood of each token, imposing local support on information transfer}
}

\newglossaryentry{neural-operator}{
    name={neural operator},
    sort={neural operator},
    description={Learned parametric mapping between function spaces rather than only between fixed-dimensional vectors}
}

\newglossaryentry{operator-learning}{
    name={operator learning},
    sort={operator learning},
    description={Machine-learning framework for approximating mappings that transform input functions, fields, or distributions into output functions, fields, or distributions}
}

\newglossaryentry{reparameterization-trick}{
    name={reparameterization trick},
    sort={reparameterization trick},
    description={Differentiable sampling construction used in variational autoencoders, typically written as a deterministic transformation of Gaussian noise}
}

\newglossaryentry{surrogate-model}{
    name={surrogate model},
    sort={surrogate model},
    description={Computationally inexpensive approximation of a more costly simulator, experiment, or physical model}
}

\newglossaryentry{transformer}{
    name={Transformer},
    sort={Transformer},
    description={Sequence-model architecture based on attention layers and token-wise transformations}
}

\newglossaryentry{under-dispersion}{
    name={under-dispersion},
    sort={under dispersion},
    description={Condition in which the variability predicted by an ensemble or probabilistic model is systematically smaller than the observed prediction error}
}

\newglossaryentry{variational-autoencoder}{
    name={variational autoencoder},
    sort={variational autoencoder},
    description={Probabilistic autoencoder that learns an approximate latent posterior and is trained using reconstruction and prior-regularization terms}
}

\newglossaryentry{conditional-variational-autoencoder}{
    name={conditional variational autoencoder},
    sort={conditional variational autoencoder},
    description={Variational autoencoder whose posterior and generation process are conditioned on auxiliary physical or machine variables}
}

\newglossaryentry{fourier-neural-operator}{
    name={Fourier Neural Operator},
    sort={Fourier Neural Operator},
    description={Neural-operator architecture that parameterizes non-local integral transformations through learned multiplication of Fourier modes}
}

\newglossaryentry{graph-neural-operator}{
    name={Graph Neural Operator},
    sort={Graph Neural Operator},
    description={Neural-operator architecture that discretizes a learned kernel through graph-based message passing on arbitrary or physically structured nodes}
}

% ============================================================================
% SYMBOLS
% ============================================================================

\newglossaryentry{sym:xi}{
    type=symbols,
    name={\ensuremath{\boldsymbol{\xi}}},
    sort={xi},
    description={Six-dimensional phase-space coordinate, typically $\boldsymbol{\xi}=(x,p_x,y,p_y,\zeta,\delta)$}
}

\newglossaryentry{sym:x}{
    type=symbols,
    name={\ensuremath{x}},
    sort={x},
    description={Horizontal transverse position}
}

\newglossaryentry{sym:px}{
    type=symbols,
    name={\ensuremath{p_x}},
    sort={px},
    description={Horizontal normalized canonical momentum}
}

\newglossaryentry{sym:y}{
    type=symbols,
    name={\ensuremath{y}},
    sort={y},
    description={Vertical transverse position}
}

\newglossaryentry{sym:py}{
    type=symbols,
    name={\ensuremath{p_y}},
    sort={py},
    description={Vertical normalized canonical momentum}
}

\newglossaryentry{sym:zeta}{
    type=symbols,
    name={\ensuremath{\zeta}},
    sort={zeta},
    description={Longitudinal position relative to the synchronous particle}
}

\newglossaryentry{sym:delta}{
    type=symbols,
    name={\ensuremath{\delta}},
    sort={delta},
    description={Relative momentum or energy deviation}
}

\newglossaryentry{sym:rho}{
    type=symbols,
    name={\ensuremath{\rho}},
    sort={rho},
    description={Beam probability density in six-dimensional phase space}
}

\newglossaryentry{sym:lambda}{
    type=symbols,
    name={\ensuremath{\lambda}},
    sort={lambda},
    description={Normalized longitudinal line density}
}

\newglossaryentry{sym:mu}{
    type=symbols,
    name={\ensuremath{\boldsymbol{\mu}}},
    sort={mu},
    description={Vector of machine, lattice, material, or collective-effect conditioning parameters}
}

\newglossaryentry{sym:phi-map}{
    type=symbols,
    name={\ensuremath{\Phi_{\boldsymbol{\mu}}}},
    sort={Phi},
    description={Single-particle or macroparticle-cloud transport map for parameters $\boldsymbol{\mu}$}
}

\newglossaryentry{sym:operator}{
    type=symbols,
    name={\ensuremath{\mathcal{G}}},
    sort={G},
    description={Target evolution or solution operator acting on a function or beam distribution}
}

\newglossaryentry{sym:empirical}{
    type=symbols,
    name={\ensuremath{\rho^{N}}},
    sort={rho N},
    description={Empirical beam measure formed from a finite macroparticle cloud}
}

\newglossaryentry{sym:centroid}{
    type=symbols,
    name={\ensuremath{m_i}},
    sort={mi},
    description={Centroid of coordinate $\xi_i$}
}

\newglossaryentry{sym:sigma}{
    type=symbols,
    name={\ensuremath{\sigma_i}},
    sort={sigma i},
    description={Coordinate-wise RMS spread of coordinate $\xi_i$}
}

\newglossaryentry{sym:emittance}{
    type=symbols,
    name={\ensuremath{\varepsilon_{x,y}}},
    sort={epsilon},
    description={Horizontal or vertical geometric emittance}
}

\newglossaryentry{sym:beta}{
    type=symbols,
    name={\ensuremath{\beta_{x,y}}},
    sort={beta},
    description={Horizontal or vertical Twiss beta function}
}

\newglossaryentry{sym:dispersion}{
    type=symbols,
    name={\ensuremath{D_{x,y}}},
    sort={D},
    description={Horizontal or vertical dispersion function}
}

\newglossaryentry{sym:eta}{
    type=symbols,
    name={\ensuremath{\eta}},
    sort={eta},
    description={Phase-slip factor in longitudinal dynamics}
}

\newglossaryentry{sym:wake-long}{
    type=symbols,
    name={\ensuremath{W_{\parallel}}},
    sort={W parallel},
    description={Longitudinal wake function}
}

\newglossaryentry{sym:wake-trans}{
    type=symbols,
    name={\ensuremath{W_{\perp}}},
    sort={W perpendicular},
    description={Transverse wake function}
}

\newglossaryentry{sym:impedance-long}{
    type=symbols,
    name={\ensuremath{Z_{\parallel}}},
    sort={Z parallel},
    description={Longitudinal beam coupling impedance}
}

\newglossaryentry{sym:impedance-trans}{
    type=symbols,
    name={\ensuremath{Z_{\perp}}},
    sort={Z perpendicular},
    description={Transverse beam coupling impedance}
}

\newglossaryentry{sym:latent}{
    type=symbols,
    name={\ensuremath{\mathbf{z}}},
    sort={z latent},
    description={Latent random variable or encoded latent state}
}

\newglossaryentry{sym:condition}{
    type=symbols,
    name={\ensuremath{\mathbf{c}}},
    sort={c},
    description={Conditioning vector supplied to a conditional model}
}

\newglossaryentry{sym:posterior}{
    type=symbols,
    name={\ensuremath{q_{\phi}(\mathbf{z}\mid X,\mathbf{c})}},
    sort={q phi},
    description={Approximate conditional latent posterior parameterized by encoder parameters $\phi$}
}

\newglossaryentry{sym:decoder}{
    type=symbols,
    name={\ensuremath{p_{\psi}(X\mid\mathbf{z},\mathbf{c})}},
    sort={p psi},
    description={Conditional decoder distribution parameterized by $\psi$}
}

\newglossaryentry{sym:histogram}{
    type=symbols,
    name={\ensuremath{H_{ij}}},
    sort={Hij},
    description={Two-dimensional histogram or probability-mass representation of coordinates $(\xi_i,\xi_j)$}
}

\newglossaryentry{sym:np}{
    type=symbols,
    name={\ensuremath{N_p}},
    sort={Np},
    description={Number of macroparticles per beam cloud}
}

\newglossaryentry{sym:ns}{
    type=symbols,
    name={\ensuremath{N_s}},
    sort={Ns},
    description={Number of simulated input--output samples or trajectories}
}

\newglossaryentry{sym:nb}{
    type=symbols,
    name={\ensuremath{N_b}},
    sort={Nb},
    description={Number of histogram bins per coordinate}
}

\newglossaryentry{sym:nzeta}{
    type=symbols,
    name={\ensuremath{N_{\zeta}}},
    sort={N zeta},
    description={Number of points in the longitudinal discretization grid}
}

\newglossaryentry{sym:r}{
    type=symbols,
    name={\ensuremath{R}},
    sort={R},
    description={Number of stochastic posterior samples used in a Monte Carlo ensemble}
}

\newglossaryentry{sym:alpha}{
    type=symbols,
    name={\ensuremath{\alpha}},
    sort={alpha},
    description={Under-relaxation coefficient in the longitudinal fixed-point iteration}
}

\newglossaryentry{sym:qfactor}{
    type=symbols,
    name={\ensuremath{Q}},
    sort={Q},
    description={Quality factor of a resonator model}
}

\newglossaryentry{sym:w1}{
    type=symbols,
    name={\ensuremath{W_1}},
    sort={W1},
    description={Wasserstein-1 distance}
}

\newglossaryentry{sym:loss}{
    type=symbols,
    name={\ensuremath{\mathcal{L}}},
    sort={L},
    description={Generic training objective or loss function}
}

% ============================================================================
% PRINTING COMMAND
% ============================================================================

% Calling this command adds every defined entry, even if the body text has not
% yet been converted to \gls{...} commands.
\newcommand{\printthesisglossaries}{%
    \glsaddall
    \printnoidxglossary[
        type=\acronymtype,
        title={List of Acronyms},
        sort=letter,
        style=long
    ]
    \printnoidxglossary[
        type=main,
        title={Glossary},
        sort=letter,
        style=long
    ]
    \printnoidxglossary[
        type=symbols,
        title={List of Symbols},
        sort=letter,
        style=long
    ]
}
